% Cote 119: offprints, reports, courses, CV and letters (typed and printed part): archive cote 119, site page 185 (PDF page 186).
% typescript. Draft compared: none.
% AI-transcribed by Claude Opus 5.5 (claude-opus-5-5, adaptive thinking, effort high), Message Batches API, batch r119_full (msgbatch_01DwZNm5Y7YsKBWyAN4sFs8w), 2026-11-13; not yet checked by eye.
\scanpage{119/185}
\mspage{- 4 -}

faisons pas l'injure de la recommencer.

8. \uline{Développement de Laurent}.

8.1. \uline{Fonctions méromorphes}.

Soient $f$ et $g$ deux fonctions analytiques dans l'ouvert connexe $U$
et supposons que $g$ ne soit pas identiquement nulle. La fonction $\frac{f(x)}{g(x)}$ est
définie et analytique au voisinage de tout point $x_o$ tel que $g(x_o) \neq 0$ c'est-%
à-dire en tout point de $U$ sauf peut-être en des points isolés (cf. \unsure{fonctinns|fonctions}
analytiques 1.5 et 2.1). Examinons le comportement de $\frac{f(x)}{g(x)}$ au voisinage d'un
point $x_o$ qui annule $g(x)$. Si $f$ n'est pas identiquement nulle on a :
\[ f(x) = (x - x_o)^k \; f_1(x) \qquad\qquad g(x) = (x - x_o)^{k'} \; g_1(x) \]
avec $k \geq 0$, $k' > 0$, $f_1$ et $g_1$ analytiques non nulles en $x_o$. $k_1(x) = \frac{f_1(x)}{g_1(x)}$
est donc analytique non nulle au voisinage de $x_o$.
Deux cas sont possibles :

1°) \uline{$k \geq k'$}

La fonction $(x - x_o)^{k-k'} \; k_1(x)$ est analytique au voisinage de $x_o$
et coïncide avec $\frac{f(x)}{g(x)}$ pour $x \neq x_o$. Ainsi prolongée en $x_o$ la fonction $\frac{f}{q}$
est analytique pour $x \neq x_o$ et a $x_o$ pour 0 si $k > k'$.

2°) \uline{$k < k'$}
\[ \frac{f(x)}{g(x)} = \frac{1}{(x-x_o)^{k'-k}} \; k_1(x) \qquad\qquad k_1(x_o) \neq 0 \]

On dit que $x_o$ est un pôle pour la fonction $\frac{f}{q}$. L'entier $k'-k$ s'appelle
l'ordre de multiplicité de ce pôle.

$\frac{f}{g}$ sera dite méromorphe sur $U$. Plus exactement :

8.1.1. \uline{Définition} :

On appelle fonction méromorphe dans un ouvert \uline{connexe} $U$ une fonc-%
tion $f(x)$ définie et analytique dans un ouvert $U'$ obtenu en enlevant de $U$
un ensemble de points isolés dont chacun est un pôle pour $f(x)$ [c'est-à-dire

\mspage{\hand{185}}
